%%%%%%%%%%%%%%%%%%%%%%%%%%%%%%%%%%%%%%%%%%%%%%%%%%%%%%%%%%%%%%%%%%%%%%%%%%%%%%%
%% Lista de Símbolos (pacote nomencl)
%%
%% Contém apenas os símbolos que aparecem nas equações dos Capítulos 3 e 4.
%% Grupos: [S] modelo de sinal, [M] modelo mecânico, [D] descritores e
%% avaliação. A macro \nomunit acrescenta a unidade alinhada à direita.
%%%%%%%%%%%%%%%%%%%%%%%%%%%%%%%%%%%%%%%%%%%%%%%%%%%%%%%%%%%%%%%%%%%%%%%%%%%%%%%

\renewcommand{\nomname}{Lista de Símbolos}

%----------------------------------------------
\newcommand{\nomunit}[1]{%
\renewcommand{\nomentryend}{\hspace*{\fill}#1}}
%----------------------------------------------
\renewcommand\nomgroup[1]{%
   \item[\bfseries
   \ifstrequal{#1}{S}{Modelo de sinal e aquisição}{%
   \ifstrequal{#1}{M}{Modelo mecânico e térmico}{%
   \ifstrequal{#1}{D}{Descritores e avaliação}{}}}%
]}

%% --- Modelo de sinal e aquisição -------------------------------------------
\nomenclature[S]{\(a(t)\)}{Aceleração simulada\nomunit{\si{\meter\per\second\squared}}}
\nomenclature[S]{\(f_r\)}{Frequência de rotação do eixo\nomunit{\si{\hertz}}}
\nomenclature[S]{\(A_k\)}{Amplitude da \(k\)-ésima harmônica da rotação\nomunit{\si{\meter\per\second\squared}}}
\nomenclature[S]{\(\varphi_k\)}{Fase da \(k\)-ésima harmônica\nomunit{\si{\radian}}}
\nomenclature[S]{\(D\)}{Amplitude do impacto do defeito\nomunit{\si{\meter\per\second\squared}}}
\nomenclature[S]{\(\zeta_d\)}{Taxa de amortecimento do impacto\nomunit{\si{\per\second}}}
\nomenclature[S]{\(f_{\mathrm{res}}\)}{Frequência da ressonância estrutural excitada\nomunit{\si{\hertz}}}
\nomenclature[S]{\(t_i\)}{Instante do \(i\)-ésimo impacto\nomunit{\si{\second}}}
\nomenclature[S]{\(\delta_i\)}{Desvio aleatório que modela o escorregamento das esferas\nomunit{\si{\second}}}
\nomenclature[S]{\(\varepsilon(t)\)}{Ruído gaussiano branco\nomunit{\si{\meter\per\second\squared}}}
\nomenclature[S]{\(u(\cdot)\)}{Degrau unitário}
\nomenclature[S]{\(f_s\)}{Taxa de amostragem\nomunit{\si{\hertz}}}
\nomenclature[S]{\(f_{\mathrm{ap}}\)}{Frequência aparente sob subamostragem sem filtro\nomunit{\si{\hertz}}}
\nomenclature[S]{\(v_{\mathrm{rms}}\)}{Velocidade eficaz de banda larga, indicador normativo\nomunit{\si{\milli\meter\per\second}}}

\nomenclature[S]{\(N_b\)}{Número de esferas do rolamento}
\nomenclature[S]{\(d\)}{Diâmetro da esfera\nomunit{\si{\milli\meter}}}
\nomenclature[S]{\(D_p\)}{Diâmetro primitivo do rolamento\nomunit{\si{\milli\meter}}}
\nomenclature[S]{\(\theta\)}{Ângulo de contato do rolamento\nomunit{\si{\radian}}}
\nomenclature[S]{\(\Delta f\)}{Resolução em frequência da DFT\nomunit{\si{\hertz}}}
\nomenclature[S]{\(T_s\)}{Duração de um símbolo LoRa\nomunit{\si{\second}}}

%% --- Modelo mecânico e térmico ---------------------------------------------
\nomenclature[M]{\(k\)}{Rigidez equivalente do caminho de transmissão\nomunit{\si{\newton\per\meter}}}
\nomenclature[M]{\(k_{\mathrm{base}}\)}{Rigidez da base, placa circular engastada\nomunit{\si{\newton\per\meter}}}
\nomenclature[M]{\(k_{\mathrm{res}}\)}{Rigidez do ressalto de centragem, compressão axial\nomunit{\si{\newton\per\meter}}}
\nomenclature[M]{\(k_{\mathrm{corpo}}\)}{Rigidez do corpo, viga tubular engastada\nomunit{\si{\newton\per\meter}}}
\nomenclature[M]{\(\mathcal{D}\)}{Rigidez flexional da placa\nomunit{\si{\newton\meter}}}
\nomenclature[M]{\(E\)}{Módulo de elasticidade\nomunit{\si{\pascal}}}
\nomenclature[M]{\(\nu\)}{Coeficiente de Poisson}
\nomenclature[M]{\(I\)}{Momento de inércia da seção tubular quadrada\nomunit{\si{\meter\tothe{4}}}}
\nomenclature[M]{\(A\)}{Área da seção resistente\nomunit{\si{\meter\squared}}}
\nomenclature[M]{\(L\)}{Comprimento do elo\nomunit{\si{\meter}}}
\nomenclature[M]{\(h\)}{Espessura da base\nomunit{\si{\meter}}}
\nomenclature[M]{\(a\)}{Raio da placa circular\nomunit{\si{\meter}}}
\nomenclature[M]{\(\ell\)}{Lado externo da seção tubular\nomunit{\si{\meter}}}
\nomenclature[M]{\(t\)}{Espessura da parede tubular\nomunit{\si{\meter}}}
\nomenclature[M]{\(f_n\)}{Frequência natural de montagem do conjunto\nomunit{\si{\hertz}}}
\nomenclature[M]{\(R_{\mathrm{th}}\)}{Resistência térmica do caminho condutivo\nomunit{\si{\kelvin\per\watt}}}

%% --- Descritores e avaliação ------------------------------------------------
\nomenclature[D]{\(x_{\mathrm{rms}}\)}{Valor eficaz do sinal na janela\nomunit{\si{\meter\per\second\squared}}}
\nomenclature[D]{\(\kappa\)}{Curtose do sinal na janela}
\nomenclature[D]{\(\sigma\)}{Desvio-padrão do sinal na janela\nomunit{\si{\meter\per\second\squared}}}
\nomenclature[D]{\(s(x,n)\)}{Escore de anomalia do Isolation Forest}
\nomenclature[D]{\(\mathrm{AUC}_\alpha\)}{Área parcial sob a curva ROC, acumulada até \(\alpha\)}
\nomenclature[D]{\(\alpha\)}{Limite superior da faixa de falso alarme na área parcial}
\nomenclature[D]{\(q\)}{Risco alvo na calibração de limiar por valores extremos}
\nomenclature[D]{\(N\)}{Número de amostras na janela}

%% --- Acrescentados pela fundamentação teórica ---------------------------------
\nomenclature[M]{\(r\)}{Razão entre frequência de excitação e frequência natural}
\nomenclature[M]{\(T(r)\)}{Transmissibilidade de base}
\nomenclature[M]{\(\zeta\)}{Razão de amortecimento}
\nomenclature[M]{\(R\)}{Resistência térmica\nomunit{\si{\kelvin\per\watt}}}
\nomenclature[D]{\(s(x,\psi)\)}{Escore de anomalia do Isolation Forest}
\nomenclature[D]{\(h(x)\)}{Comprimento do caminho de isolamento de \(x\)}
\nomenclature[D]{\(c(\psi)\)}{Comprimento médio de caminho para subamostra de \(\psi\) pontos}
\nomenclature[D]{\(\psi\)}{Tamanho da subamostra de cada árvore}
\nomenclature[D]{\(\beta\)}{Limite superior de falso alarme da área parcial}
